%%%%%%%%%%%%%%%%%%%%%%%%
% $Autor: MansourAhmadyPhoulady, RanjeetsingTawar, HemanthPrabhakaran$
% $Pfad: Car_Licence_Plate_Identification_with_a_JetsonNano$
% $Version: 4.0.0 $
%
% !TeX encoding = utf8
%
%%%%%%%%%%%%%%%%%%%%%%%%


\chapter{Requirements.txt}

The \texttt{requirements.txt} file serves as a crucial component of Python projects, as it contains a comprehensive list of all the Python packages and their specific versions required to successfully run the project. This file ensures that the project dependencies are consistently installed across different environments, thereby enhancing reproducibility and minimizing potential compatibility issues.

Below is the detailed list of the packages and their respective versions required for this project:

\begin{verbatim}
	# requirements.txt
	
	# List of Python packages and their versions for the project dependencies.
	
	absl-py==2.1.0
	apache-beam==2.54.0rc1
	asttokens==2.4.1
	astunparse==1.6.3
	avro-python3==1.10.2
	cachetools==5.3.2
	certifi==2024.2.2
	charset-normalizer==3.3.2
	colorama==0.4.6
	comm==0.2.1
	contextlib2==21.6.0
	contourpy==1.2.0
	cycler==0.12.1
	Cython==3.0.8
	debugpy==1.8.0
	decorator==5.1.1
	dm-tree==0.1.8
	executing==2.0.1
	flatbuffers==23.5.26
	fonttools==4.47.2
	gast==0.5.4
	google-api-core==2.8.0
	google-auth==2.27.0
	google-auth-oauthlib==1.2.0
	google-cloud-core==2.3.0
	google-cloud-storage==2.3.0
	google-crc32c==1.5.0
	google-pasta==0.2.0
	google-resumable-media==2.3.3
	googleapis-common-protos==1.56.1
	grpcio==1.60.1
	grpcio-status==1.33.2
	h5py==3.10.0
	idna==3.6
	ipykernel==6.29.0
	ipython==8.21.0
	jedi==0.19.1
	jupyter-client==8.6.0
	jupyter-core==5.7.1
	keras==2.15.0
	kiwisolver==1.4.5
	libclang==16.0.6
	lvis==0.5.3
	lxml==5.1.0
	Markdown==3.5.2
	MarkupSafe==2.1.5
	matplotlib==3.8.2
	matplotlib-inline==0.1.6
	ml-dtypes==0.2.0
	namex==0.0.7
	nest-asyncio==1.6.0
	numpy==1.26.3
	oauthlib==3.2.2
	object-detection==0.1
	opencv-python==4.9.0.80
	opt-einsum==3.3.0
	packaging==23.2
	pandas==2.2.0
	parso==0.8.3
	pillow==10.2.0
	pip==24.0
	platformdirs==4.2.0
	portalocker==2.8.2
	prompt-toolkit==3.0.43
	protobuf==3.20.0
	psutil==5.9.8
	pure-eval==0.2.2
	pyasn1==0.5.1
	pyasn1-modules==0.3.0
	pycocotools==2.0.7
	Pygments==2.17.2
	pyparsing==2.4.7
	python-dateutil==2.8.2
	pytz==2024.1
	pywin32==306
	PyYAML==6.0.1
	pyzmq==25.1.2
	regex==2023.12.25
	requests==2.31.0
	requests-oauthlib==1.3.1
	rich==13.7.0
	rsa==4.9
	sacrebleu==2.2.0
	scipy==1.12.0
	setuptools==65.5.0
	six==1.16.0
	slim==0.1
	stack-data==0.6.3
	tabulate==0.9.0
	tensorboard==2.15.1
	tensorboard-data-server==0.7.2
	tensorflow==2.15.0
	tensorflow-estimator==2.15.0
	tensorflow-intel==2.15.0
	tensorflow-io==0.31.0
	tensorflow-io-gcs-filesystem==0.31.0
	termcolor==2.4.0
	tf-models-official==2.15.0
	tf-slim==1.1.0
	tornado==6.4
	traitlets==5.14.1
	typing_extensions==4.9.0
	urllib3==2.2.0
	wcwidth==0.2.13
	Werkzeug==3.0.1
	wget==3.2
	wheel==0.42.0
\end{verbatim}

\section{Steps to Create a requirements.txt File}

The \texttt{requirements.txt} file can be generated through several methods. Here are the commonly used approaches:

\begin{enumerate}
	\item \textbf{Using \texttt{pip freeze}:} To generate the file directly from an environment, run the following command in your terminal:
	\begin{lstlisting}[language=bash]
		pip freeze > requirements.txt
	\end{lstlisting}
	This command lists all installed packages in the current Python environment along with their versions and saves them into a file named \texttt{requirements.txt}.
	
	\item \textbf{Manually Creating the File:} If you already know the packages required for your project, you can manually create a \texttt{requirements.txt} file. Each package name should be listed on a new line, optionally followed by a specific version number for consistency across environments.
	
	\item \textbf{From a Virtual Environment:} Using a virtual environment is a best practice to manage dependencies in Python projects. After activating the virtual environment and installing the necessary packages, use the \texttt{pip freeze} command to export the installed dependencies into a \texttt{requirements.txt} file. This ensures isolation and prevents conflicts between dependencies in different projects.
\end{enumerate}

\section{Benefits of Using requirements.txt}

\begin{itemize}
	\item \textbf{Reproducibility:} Ensures that the same versions of dependencies are installed across different environments, minimizing discrepancies during project execution.
	\item \textbf{Ease of Deployment:} Simplifies the process of setting up new environments by automating the installation of required dependencies.
	\item \textbf{Collaboration:} Facilitates smooth collaboration among team members by ensuring consistency in the development environment.
	\item \textbf{Version Control:} Helps in tracking and managing changes to dependencies over time, which is crucial for maintaining and upgrading projects.
\end{itemize}

By leveraging a \texttt{requirements.txt} file, developers can streamline the development, deployment, and maintenance processes, ensuring smooth operation and efficient collaboration within teams.
